\documentclass[10pt,a4paper]{amsart}
\usepackage[margin=19mm]{geometry}
\usepackage{amsmath,amssymb,mathtools}
\usepackage[scr=rsfs,cal=euler]{mathalfa}
\usepackage{xfrac}
\usepackage[unicode,hidelinks,bookmarks=false]{hyperref}
\newcommand{\ie}{i.e.\ }
\newcommand{\cf}{cf.\ }
\newcommand{\resp}{respectively\ }
\newcommand{\Subsection}{\subsection}
\providecommand{\nobreakdash}{\nobreak\hbox{-}}
\setlength{\emergencystretch}{2em}
\multlinegap0pt
\usepackage{amssymb}
\usepackage{mathtools}
\usepackage[shortlabels]{enumitem}
\newtheorem{lemma}{Lemma}
\newtheorem{proposition}{Proposition}
\title{An Improved Lower Bound for the de Bruijn--Erd\H{o}s Consecutive Gap Problem}
\author{Samuel Korsky}
\date{Source: arXiv:2605.30959v1 (CC BY 4.0)}
\begin{document}
\maketitle

\noindent\emph{Source and licence.} An Improved Lower Bound for the de Bruijn--Erd\H{o}s Consecutive Gap Problem. Version arXiv:2605.30959v1 (CC BY 4.0), distributed by arXiv under the Creative Commons Attribution 4.0 International licence, http://creativecommons.org/licenses/by/4.0/, as recorded in the arXiv licence metadata for this exact version and shown on the abstract page https://arxiv.org/abs/2605.30959v1. Full licence notice: \texttt{licenses/arxiv-2605.30959-CC-BY-4.0.txt}.\par\smallskip

\noindent\textit{Editorial scope.} Counted unit: the article's proposition prop:epoch (source0.tex lines 319--324). The article's complete original proof of it is delivered with it (source0.tex lines 326--390, one \texttt{proof} environment of 65 lines). No mathematics is written, repaired or re-derived: the statement and the proof are the article's own lines, in the article's own notation, and no retained line is substituted or re-flowed.  Retained uncounted as the source-local context the proof needs: lines 167--212.  Nothing was omitted from any retained span, so the package carries no omission record.  No retained line contains a citation, so the wrapper carries no bibliography and no bibliography entry.  No retained line contains a figure environment, an image-inclusion command or drawing code, so no image is delivered and the package declares no asset dependency.  Editorial formatting only, in addition: an AMS \texttt{amsart} preamble that restates the article's own declarations and macros used by the retained lines, namely the environments \texttt{lemma}, \texttt{proposition} on separate counters, together with the standard packages those lines need.  The retained environments print in this document as Lemma 1, Proposition 1 in order of appearance, whereas the article prints the same items with the numbering of its own full document.  The complete native source of the article as distributed in the arXiv e-print for this exact version is bundled unchanged as \texttt{sources/201704/source0.tex}.
\begin{lemma}[Protected block]\label{lem:protected-block}
Fix $\eta\in(0,1)$ and suppose that the step $n\to n+1$ satisfies
\[
    M_{n+1}\geq (1-\eta)M_n
\]
and
\[
    R_{n+1}\leq \rho.
\]
Suppose the split gap is $\ell=x+y$. After the split, consider the $2r$ consecutive gaps
\[
    h_1,h_2,\ldots,h_{2r},
\]
where
\[
    h_r=x,\qquad h_{r+1}=y,
\]
with $r-1$ gaps to the left of $\ell$ and $r-1$ gaps to the right of $\ell$ included.

\bigskip
\noindent
Set
\[
    q=\frac{1-\eta}{\rho},
    \qquad
    \alpha=1-q=\frac{\rho-1+\eta}{\rho}.
\]
Then
\[
    h_j\leq \alpha M_n
    \qquad (1\leq j\leq 2r).      \tag{3.1}\label{eq:gap-bound}
\]

\smallskip
\noindent
Moreover, suppose that at a later time $T$ one of the gaps $h_j$ is split, and that no gap among
\[
    h_1,\ldots,h_{2r}
\]
has been split before time $T$. If $R_T\leq \rho$, then
\[
    M_T\leq \beta M_n,
    \qquad
    \beta=(r-1)(\rho-1+\eta).       \tag{3.2}\label{eq:release-factor}
\]
\end{lemma}

\begin{proposition}\label{prop:epoch}
There is a constant $C=C(r,\eta,\rho)$ such that, for all sufficiently large $N$,
\[
    N^+\leq \left(1+\frac1r\right)N+C .
\]
\end{proposition}

\begin{proof}
Before time $N^+$, we have
\[
    M_n>\beta M_N .
\]
Consider a slow step $k\to k+1$ with $N\leq k<N^+-1$. It creates a marked block as in Lemma \ref{lem:protected-block}, with $M_k\leq M_N$. By the lemma, none of the marked gaps can be split at any later time $T<N^+$, since otherwise
\[
    M_T\leq \beta M_k\leq \beta M_N,
\]
contradicting the definition of $N^+$. Thus every slow split before the terminal step $N^+-1\to N^+$ creates a protected block that remains untouched until time $N^+$.

\bigskip
\noindent
First we bound the number of fast steps before the terminal step. Suppose there are $b$ fast steps among
\[
    N\to N+1\to\dots\to N^+-2\to N^+-1.
\]
At every fast step $M_n$ is multiplied by a factor less than $1-\eta$, while at all other steps it does not increase. Therefore
\[
    M_{N^+-1}\leq (1-\eta)^b M_N.
\]
But by minimality of $N^+$,
\[
    M_{N^+-1}>\beta M_N.
\]
Consequently
\[
    (1-\eta)^b>\beta,
\]
so
\[
    b\leq B_0:=\left\lceil\frac{\log \beta}{\log(1-\eta)}\right\rceil .
\]
Here $B_0$ depends only on $r,\eta,\rho$.

\bigskip
\noindent
It remains to count the slow steps. We compare them with the $N$ gaps present at time $N$. Let $F$ be the set of initial gaps whose descendants are split during a fast step before the terminal step. Then $|F|\leq B_0$. Call an initial gap ``bad" if its cyclic distance, in the initial $N$-cycle, from some gap in $F$ is at most $r$. There are $O_r(B_0)$ bad initial gaps.

\bigskip
\noindent
The slow steps associated to bad initial gaps contribute only $O_{r,\eta,\rho}(1)$. Indeed, initially there are only $O_r(B_0)$ bad gaps. During the epoch, a fast split can increase by at most one the number of active descendants of bad gaps, and there are at most $B_0$ fast splits. A slow split of an active descendant removes it from further consideration, since its children lie in the protected block created by that slow split and cannot be split again before $N^+$. Thus the number of slow splits associated to bad initial gaps is $O_r(B_0)+B_0=O_{r,\eta,\rho}(1)$.

\bigskip
\noindent
For each remaining slow split, associate it to the unique initial gap whose descendant is split. We claim that these associated initial gaps are pairwise at cyclic distance at least $r$. Suppose not. Then two remaining slow splits are associated to good initial gaps $I$ and $J$ whose cyclic distance is at most $r-1$, allowing $I=J$. Let $A$ be the shorter cyclic arc of initial gaps from $I$ to $J$. Then $A$ contains at most $r$ initial gaps. Since $I$ and $J$ are good, no gap in $A$ has a descendant split during a fast step before the terminal step $N^+-1\to N^+$.

\bigskip
\noindent
Let $k\to k+1$ be the first slow split before the terminal step associated to an initial gap in $A$. Since no gap in $A$ has previously been split by either a fast or a slow step, all gaps of $A$ are still intact immediately before time $k$. At least one of the two chosen slow splits occurs after time $k$; let $J'\in A$ be the initial gap associated to such a later split. At time $k$, the gap $J'$ either is the split gap itself or lies among the $r-1$ gaps to the left or the $r-1$ gaps to the right of the split gap. Hence, after the split at time $k$, the descendant of $J'$ lies inside the protected block created at time $k$. This descendant therefore cannot be split before $N^+$, a contradiction.

\bigskip
\noindent
Thus, apart from $O_{r,\eta,\rho}(1)$ discarded exceptions, the initial gaps associated to slow split locations form a subset of an $N$-cycle with mutual cyclic distance at least $r$. Such a subset has size at most $N/r$. Therefore the number of slow steps before the terminal step is at most
\[
    \frac Nr+O_{r,\eta,\rho}(1).
\]
Adding the bounded number of fast steps and the terminal step $N^+-1\to N^+$ gives
\[
    N^+-N
    \leq
    \frac Nr+O_{r,\eta,\rho}(1),
\]
which proves the proposition.
\end{proof}

\end{document}
